\cotefile{81}
\scanpage{27}
% loc: 1–3
cette dernière relation $\Sigma_1^-\cap\Sigma_2^-\ne\varnothing$ est vraie \unsure{=|même} si $E_1=E_2=\varnothing$, car \deleted{\unsure{alors par hyp.|par hyp.}} si on avait de plus $\Sigma_1^-\cap\Sigma_2^-=\varnothing$, on en
% loc: 4–8
conclurait que les deux \deleted{\unsure{part.}} décompositions binaires \deleted{\unsure{son|se}} (qui sont alors des partitions binaires) sont \uline{identiques}, ce que l'on \inserted{\unsure{a|à}} exclu. Cela montre que les $\Sigma_1^+,\Sigma_2^+$ sont uniquement déterminés.

% loc: 9
On a donc le diagramme d'\unsure{inclusions|inclusion}
% loc: upper chain
\diagram{\begin{tikzcd}[ampersand replacement=\&,column sep=small]
\Sigma_1^+ \& E_1 \arrow[l,hook] \arrow[r,hook] \& \Sigma_{12} \& E_2 \arrow[l,hook] \arrow[r,hook] \& \Sigma_2^+
\end{tikzcd}}
% loc: 10
\unsure{qu'on|qui est|qu'on a}, \unsure{exprime|exprimé} par \unsure{nos diagrammes|un diagramme}.
% loc: lower chain
\diagram{\begin{tikzcd}[ampersand replacement=\&,column sep=small,row sep=2pt]
\unsure{\bullet|\circ} \arrow[rr,dash] \& {} \& \unsure{\bullet|\blacktriangledown} \arrow[rr,dash] \& {} \& \bullet\\
\Sigma_1^+ \& E_1 \arrow[l,hook] \arrow[r,hook] \& \Sigma_{12} \& E_2 \arrow[l,hook] \arrow[r,hook] \& \Sigma_2^+
\end{tikzcd}}
% loc: NB at right of lower chain
\margin{\[
\text{NB}\left\{\begin{aligned}
\Sigma_1^-&=\Sigma_{12}\amalg_{E_2}\Sigma_2^+\\
\Sigma_{\deleted{\unsure{1}}2}^-&=\deleted{\unsure{E}}\Sigma_{12}\amalg_{E_1}\Sigma_1^+
\end{aligned}\right.
\]}
% slanted
% loc: slanted note, lines 1–8
\margin{Ceci n'implique pas que les quatre ensembles $\Sigma_1^+,\Sigma_1^-,\Sigma_2^+,\Sigma_2^-$ soient distincts, car on peut avoir (\unsure{non pas|bien que|bien pas} $\Sigma_1^+=\Sigma_2^+$, ni $\Sigma_1^-=\Sigma_2^-$, mais) $\Sigma_1^+=\Sigma_2^-$, qui équivaut à $E_1=\Sigma_{12}$ et entraîne $E_2=\varnothing$
% loc: slanted note, lines 9–13
(Cela signifie, en termes des $(\Sigma_{\unsure{1|t|i}}^+,\Sigma_1^-)$ que $\Sigma_2^{\unsure{\dot{-}|-}}=\Sigma_1^+$, $\Sigma_{\unsure{2|t}}^+=\unsure{X|\varnothing}\setminus\Sigma_1^+=\overset{\unsure{\circ|\leftarrow}}{\Sigma}{}_1^{-,\unsure{\circ|0}}$) \deleted{\unsure{ou}} et la relation \unsure{symetrique|symétrique}. On ne veut \unsure{pas exclure} \unsure{ce|ces} cas, \unsure{i.e.|ici} \deleted{\unsure{mais|nous}} n'\unsure{exige pas|exigeons pas} que
% loc: slanted note, lines 14–19
les quatre ensembles $\Sigma_1^+,\Sigma_1^-,\Sigma_2^+,\Sigma_2^-$ soient distincts \unsure{s'ils|ssi ils} le \unsure{\deleted{sont}|sont}, \unsure{on|c.-à-d.} \unsure{dit|dire} que les \deleted{\unsure{ens.}} \inserted{deux \unsure{décomp. bin.|décomp. bis}} \deleted{\unsure{sont strictmt}} \unsure{non parallèles|sont non parallèles}.}
% loc: 11–12, main text below chains
Je me propose d'étudier \unsure{les|des} ensembles finis \unsure{indicellés|indiciellés|indécollés|indécollt} $\deleted{\unsure{\Pi|\Sigma}}(\pi_\alpha)_{\alpha\in A}\deleted{\unsure{=}(\unsure{\Sigma_\alpha^i}),\unsure{\alpha\in A}}$
% loc: 13–14
$\deleted{\pi_\alpha=\{\Sigma_\alpha^\varepsilon}$ \unsure{de|des} décompositions binaires de $\unsure{X}$ deux à deux non parallèles.

% loc: 15–16
Voici une façon \unsure{très} \unsure{générale|génerale} d'\deleted{\unsure{en}}obtenir des \unsure{floppées|floppés}.
% archive sheet number: 26
